\providecommand{\BM}{\textbf{[B]}}
\providecommand{\KM}{\textbf{[K]}}
\providecommand{\SM}{\textbf{[S]}}
\providecommand{\XM}{\textbf{[X]}}
\providecommand{\QM}{\textbf{[Q]}}

\chapter*{Gemeinsamer Anker: Brückengleichungen zu Observer Patch Holography und die elektroschwache Skala als Treffpunkt}
\addcontentsline{toc}{chapter}{Gemeinsamer Anker: Brückengleichungen zu OPH und die elektroschwache Skala als Treffpunkt}

{\small\noindent\textbf{Zusammenfassung.}\quad\ignorespaces
Observer Patch Holography (OPH) führt als offenen Punkt, dass Länge und Takt
nicht aus den eigenen Aktualisierungsregeln folgen. Dieses Dokument schreibt
die Brücke zu FFGFT als Gleichungen aus. Die Zellenenergie von OPH ist exakt
die FFGFT-Bitenergie $E_{\rm bit}=\hbar c/L$ an der Zellkante
$L_{\rm cell}=\sqrt P\,\ell_P$ \BM; die Kollapsschwelle aus Dok.~329 ergibt dort
$n_{\rm thr}=\sqrt{P/2}$ mit der Grenze $P=2$ \BM; $\tilde T\cdot m=1$ liefert
einen eigenen Takt $\tilde T_{\rm cell}=\sqrt P\,t_P$. Unterhalb von $v$ liefert FFGFT
alle Massen als Verhältnisse zu $v$. Die elektroschwache Skala $v$ ist damit der
Treffpunkt: die bedingte OPH-Hierarchie $v/E_\star$ \QM{} und die FFGFT-Leiter
$m_e/v$ \KM{} ergeben zusammen $m_e/E_P$ auf $0{,}93\,\%$. Ein einziger
SI-Anker legt dann beide Rahmen fest; über die eigene FFGFT-Kette werden
Zellkante und Zellentakt absolut. Der Koeffizient dieser Lesart ist auf
OPH-Seite inzwischen abgeleitet (Dok.~376); \#740 selbst bleibt dort offen.
}\medskip

\section*{Zeichenerklärung}
\addcontentsline{toc}{section}{Zeichenerklärung}

\medskip
\begin{tabular}{@{}ll@{}}
\textbf{[B]} & algebraisch bewiesen \\
\textbf{[K]} & numerisch bestätigt \\
\textbf{[S]} & strukturell plausibel, Herleitung offen \\
\textbf{[X]} & geprüft und negativ \\
\textbf{[Q]} & aus externer Quelle übernommen, zitiert \\
\end{tabular}

\medskip
\begin{tabular}{@{}lp{10cm}@{}}
\toprule
\textbf{Symbol} & \textbf{Bedeutung} \\
\midrule
$\xi$ & FFGFT-Skalierungsparameter ($4/30000$) \\
$\ell_P$, $t_P$, $E_P$ & Planck-Länge, -Zeit, -Energie (nicht reduziert) \\
$L_0=\xi\ell_P$ & Auflösungsboden der FFGFT (Dok.~180) \\
$E_{\rm bit}(L)=\hbar c/L$ & systemabhängige Bitenergie (Dok.~257) \\
$n_{\rm thr}(L)$ & Kollapsschwelle, Zahl der Bitquanten bis Compton $=$ Schwarzschild (Dok.~329) \\
$M_{\rm coll}=m_P/\sqrt2$ & Kollapsmasse (Dok.~329) \\
$\tilde T$ & intrinsische Zeitskala, $\tilde T\cdot m=1$ \\
$v$ & elektroschwache Skala ($246{,}22$\,GeV, aus $G_F$) \\
$r_i$, $p_i$ & Vorfaktor und Exponent der Leiter $m_i=r_i\,\xi^{p_i}v$ \\
\midrule
$P$ & OPH-Pixelverhältnis $a_{\rm cell}/\ell_P^2$ \\
$L_{\rm cell}$, $E_{\rm cell}$ & OPH-Zellkante $\sqrt P\,\ell_P$, Zellenergie $E_P/\sqrt P$ \\
$E_\star$ & OPH-Referenzenergie der Hierarchie $v/E_\star$ \\
$\alpha_U$ & OPH-Vereinigungskopplung \\
$N$ & OPH-Kapazität des Gesamtnetzes \\
\bottomrule
\end{tabular}

\section{Anlass und Einordnung}

Im Tracking-Eintrag \#740 des OPH-Repositoriums wird festgehalten, dass Länge
und Takt nicht aus den OPH-Aktualisierungen folgen: die Länge ist nur über die
Kapazität $N$ und $\Lambda\ell_P^2=3\pi/N$ an einen äußeren Wert gebunden, der
Takt nur über Verhältnisse. Dok.~364 hat den ikosaedrischen Träger von OPH
spektral untersucht; dieses Dokument behandelt die Skalenfrage.

Die Rollen sind nicht symmetrisch. FFGFT leitet aus $\xi$ ab; Messwerte und
Standardmodell-Größen dienen nur der Umrechnung und dem Vergleich. OPH-Größen
treten hier als zitierte Eingänge eines Fremdrahmens auf \QM{} und tragen den
Status, den ihre eigene Dokumentation ihnen gibt. Der kosmische Sektor bleibt
ausgeklammert (P39); $N$ wird nicht verwendet.

\section{Eingänge}

\subsection*{FFGFT}
\begin{itemize}\itemsep2pt
\item Bitenergie $E_{\rm bit}(L)=\hbar c/L$, systemabhängig (Dok.~257).
\item Kollapsschwelle $n_{\rm thr}(L)=L/(\sqrt2\,\ell_P)$ und
  $M_{\rm coll}=m_P/\sqrt2$ aus Compton $=$ Schwarzschild (Dok.~329) \BM.
\item Auflösungsboden $L_0=\xi\,\ell_P$ als Verhältnis (Dok.~180, 329).
\item $\tilde T\cdot m=1$: jeder Energie gehört eine Eigenzeit $\hbar/E$.
\item Leptonleiter $m_e=\tfrac43\,\xi^{3/2}\,v$ (Dok.~006, 373) \KM.
\end{itemize}

\subsection*{OPH \QM}
\begin{itemize}\itemsep2pt
\item Pixelverhältnis $P=a_{\rm cell}/\ell_P^2$ als Fixpunkt von
  $P=\varphi+\sqrt\pi/A_T(P)$; öffentliche Wurzel $P=1{,}630968\ldots$,
  Quellwurzel $1{,}630972\ldots$; Zellenergie $E_{\rm cell}=E_P/\sqrt P$
  (\texttt{code/P\_derivation/FULL\_DERIVATION.md}). Die öffentliche Wurzel
  enthält einen empirischen Hadronenbeitrag; die Quellwurzel nicht.
\item Hierarchie $v/E_\star=P^{-1/2}\exp[-2\pi/(4\alpha_U(P))]
  =2{,}01998\times10^{-17}$ (Theorem~A,
  \texttt{code/particles/hierarchy/theorem\_package.md}); ausdrücklich
  bedingt, die physikalische Identität von $E_\star$ ist in OPH offen.
\end{itemize}

\section{Brückengleichungen für Länge und Takt}\label{sec:bruecke}

Alle Aussagen dieses Abschnitts sind dimensionslos in Planck-Einheiten; kein
SI-Anker wird verbraucht.

\paragraph{B1 --- Wörterbuch \BM.} Mit $L_{\rm cell}=\sqrt P\,\ell_P=1{,}2771\,\ell_P$ gilt
\begin{equation}
  E_{\rm cell}^{\rm OPH}=\frac{E_P}{\sqrt P}=\frac{\hbar c}{L_{\rm cell}}=E_{\rm bit}(L_{\rm cell}).
\end{equation}
Beide Rahmen benutzen dieselbe Energie-Längen-Beziehung. OPH legt die Länge
durch seinen Fixpunkt fest; in FFGFT bleibt sie systemabhängig.

\paragraph{B2 --- Kollapsschwelle \BM.} Eingesetzt in Dok.~329:
\begin{equation}
  n_{\rm thr}(L_{\rm cell})=\sqrt{P/2}=0{,}9030,\qquad
  \frac{r_s(E_{\rm cell})}{L_{\rm cell}}=\frac2P=1{,}2263 .
\end{equation}
Die Grenze, ab der ein einzelnes Zellquant außerhalb seines eigenen
Schwarzschild-Radius bleibt, liegt bei $P=2$. Die OPH-Wurzel liegt darunter:
ein auf eine Zelle lokalisiertes Zellquant sitzt an der Kollapsskala. Für ein
holografisches Pixel ist das zu erwarten, nicht einzuwenden; es ist aber eine
scharfe Bedingung, die jede künftige Änderung von $P$ einhalten oder erklären
muss. Quell- und öffentliche Wurzel unterscheiden sich hier um
$1{,}2\times10^{-6}$; die Aussage hängt nicht am offenen Hadronenteil.

\paragraph{B3 --- Takt.} $\tilde T\cdot m=1$ ordnet der Zelle eine Eigenzeit zu:
\begin{equation}
  \tilde T_{\rm cell}=\frac{\hbar}{E_{\rm cell}}=\sqrt P\,t_P=\frac{L_{\rm cell}}{c}.
\end{equation}
Als Rechnung ist das \BM; physikalisch gibt es OPH einen Takt pro Zelle in
derselben Einheit wie die Länge, ohne Mastertakt. Die Anbindung an eine
Laboruhr ist möglich, aber auf OPH-Seite noch nicht ausgeführt \SM. Eine
Laboruhr, etwa der Cäsium-Übergang, ist selbst ein Materiesystem mit Energie
$\Delta E_{\rm clk}$ und nach $\tilde T\cdot m=1$ eigenem Takt $\hbar/\Delta E_{\rm clk}$.
Die Anbindung verlangt das Verhältnis $\Delta E_{\rm clk}/E_{\rm cell}$, aus OPH
selbst berechnet; OPH führt diesen Schritt als offen (Uhrengrenze
$\varepsilon_{\rm clk}=\Delta E_{\rm clk}/E_\star$, Cäsium-Weg als unvollständiger
Kandidat). In FFGFT ist die Frage nicht offen: jede Mode hat ihren Takt
$\tilde T=1/m$, Taktverhältnisse sind Massenverhältnisse. FFGFT liefert damit
das Übersetzungsprinzip, nicht den fehlenden OPH-Wert.

\paragraph{B4 --- Boden \KM.} $L_{\rm cell}/L_0=\sqrt P/\xi=9578$. Die
OPH-Zelle liegt rund vier Größenordnungen über dem FFGFT-Auflösungsboden; kein
Widerspruch.

\section{Der Treffpunkt $v$}\label{sec:treffpunkt}

Die beiden Rahmen decken verschiedene Abschnitte der Skalenleiter ab. OPH
rechnet bedingt von oben bis $v$, FFGFT von $v$ nach unten: alle geladenen
Leptonmassen über die Leiter $m_i=r_i\,\xi^{p_i}v$ (Dok.~006, 373), die
Bosonmassen über $\sin^2\theta_W=2/9$ und die Spurregel (Dok.~372, 373).

\begin{center}
\begin{tabular}{@{}llr@{}}
\toprule
\textbf{Größe} & \textbf{Herkunft} & \textbf{gegen Messwert} \\
\midrule
$v/E_\star=2{,}01998\times10^{-17}$ & OPH, Theorem A \QM & $+0{,}16\,\%$ (mit $E_\star=E_P$) \\
$m_e/v=\tfrac43\,\xi^{3/2}=2{,}0528\times10^{-6}$ & FFGFT, Dok.~006/373 \KM & $-1{,}09\,\%$ (bare) \\
$m_e/E_P=4{,}1466\times10^{-23}$ & Produkt & $-0{,}93\,\%$ \\
\bottomrule
\end{tabular}
\end{center}

Zur Referenzenergie $E_\star$ siehe den folgenden Unterabschnitt. Die $-1{,}09\,\%$ auf FFGFT-Seite sind der bekannte
Abstand bare zu gemessen der Leptonmassen (Dok.~352), kein neuer Befund. Das
Produkt hängt vom OPH-Zweig nur auf $10^{-4}$ ab.

Folge: Ein einziger SI-Anker legt beide Rahmen fest. Mit $G$ als Anker führt
die OPH-Hierarchie auf $v$ und die FFGFT-Leiter von dort auf alle
Teilchenmassen; mit $m_e$ als Anker läuft die Kette umgekehrt. Die
Restabweichung von rund $1\,\%$ ist der gemeinsame Konsistenztest. FFGFT erreicht
$m_e/E_P$ auch über die eigene Kette; der Treffpunkt $v$ ist damit ein
Vergleich zweier unabhängiger Wege.

\subsection*{Die Referenzenergie $E_\star$}\label{sec:estar}

\paragraph{Was $E_\star$ ist.} In OPH ist $E_\star$ die Energieeinheit der
Rechnung. Der Code setzt die Planck-Energie gleich eins; die Zellenenergie ist
$E_P/\sqrt P$, die Vereinigungsskala $M_U=E_P\,e^{-2\pi}P^{1/6}$, die Hierarchie
$v/E_\star=P^{-1/2}\exp[-2\pi/(4\alpha_U)]$ (reproduziert, Prüfskript D1). Welche
physikalische Energie diese Einheit ist, legt OPH ausdrücklich nicht fest; die
Festlegung soll über eine Uhr erfolgen, $E_\star=h\nu_{\rm clk}/\varepsilon_{\rm clk}$,
also über denselben offenen Schritt wie in B3.

\paragraph{Zwei Kandidaten.} Die Planck-Energie gibt es in zwei Konventionen:
nicht reduziert $E_P=\sqrt{\hbar c^5/G}=1{,}2209\times10^{19}$\,GeV und
reduziert $\bar E_P=E_P/\sqrt{8\pi}=2{,}435\times10^{18}$\,GeV. Sie
unterscheiden sich um $\sqrt{8\pi}\approx5{,}01$, je nachdem, ob die
Gravitationskopplung als $G$ oder als $8\pi G$ geschrieben wird. Mit gemessenem
$v$ ergibt OPH $E_\star=1{,}2189\times10^{19}$\,GeV: $0{,}16\,\%$ unter $E_P$,
aber um den Faktor $5{,}0$ über $\bar E_P$ \KM. Die reduzierte Lesart ist damit
ausgeschlossen, die nicht reduzierte passt.

\paragraph{Wie scharf $0{,}16\,\%$ ist.} $\alpha_U$ steht im Exponenten. Die
Empfindlichkeit ist
\begin{equation}
  \frac{\partial\ln(v/E_\star)}{\partial\alpha_U}=\frac{\pi}{2\alpha_U^2}\approx929 .
\end{equation}
Die Abweichung von $0{,}16\,\%$ entspricht einer Verschiebung von $\alpha_U$ um
$1{,}7\times10^{-6}$, relativ $4\times10^{-5}$ \BM. Dieselbe Empfindlichkeit
erklärt exakt den Unterschied zwischen öffentlichem und Quellzweig
($-8{,}4\times10^{-5}$, Prüfskript D5). Die Übereinstimmung mit $E_P$ ist
damit keine grobe Größenordnungsnähe, sondern eine enge Bedingung an $\alpha_U$.

\paragraph{Warum Lesart \SM.} OPH selbst behauptet $E_\star=E_P$ nicht; die
Zuordnung folgt hier aus dem Zahlenvergleich. Bestätigt wäre sie, wenn OPH die
Uhrenanbindung schließt und dabei $E_\star=E_P$ herauskommt, oder wenn in OPH
die Kopplung ohne $8\pi$ strukturell ausgezeichnet wird. Für den Vergleich in
§\ref{sec:wege} ist die Lesart unkritisch: sie legt nur fest, gegen welche der
beiden Planck-Energien die $0{,}16\,\%$ gemessen werden.

\section{Zwei Wege zu $m_e/E_P$ und ihre Genauigkeit}\label{sec:wege}

FFGFT erreicht die Planck-Skala über die eigene Kette (Dok.~012, 013, 180):

\begin{enumerate}\itemsep2pt
\item $G=\dfrac{\xi^2}{4m_e}\,C_{\rm conv}\,K_{\rm frak}=6{,}6745\times10^{-11}\,
  \mathrm{m^3kg^{-1}s^{-2}}$, gegen CODATA $+0{,}003\,\%$. Die Schreibweise in
  Dok.~180 teilt denselben Wert in $C_{\rm dim}\cdot C_{\rm conv}$ auf; beide
  Formen sind gleichwertig.
\item $\ell_P=\sqrt{\hbar G/c^3}=1{,}6163\times10^{-35}$\,m und damit
  $L_0=\xi\,\ell_P=2{,}155\times10^{-39}$\,m.
\item $E_P=1{,}22087\times10^{22}$\,MeV, $m_e/E_P=4{,}1855\times10^{-23}$
  ($+0{,}002\,\%$) \KM.
\item $v/E_P=2{,}01676\times10^{-17}$.
\end{enumerate}

Nicht jeder Rechenweg ergibt dieselbe Genauigkeit:

\begin{center}
\small
\begin{tabular}{@{}llrl@{}}
\toprule
\textbf{Weg} & \textbf{Größe} & \textbf{Abw.} & \textbf{Herkunft der Abweichung} \\
\midrule
FFGFT-Kette $\xi\to G\to E_P$ & $m_e/E_P$ & $+0{,}002\,\%$ & Genauigkeit von $G$ \\
FFGFT-Leiter & $m_e/v$ & $-1{,}09\,\%$ & bare gegen gemessen (Dok.~352) \\
OPH, Theorem A \QM & $v/E_\star$ gegen FFGFT-Kette & $+0{,}16\,\%$ & Lesart $E_\star=E_P$, bedingter Fixpunkt \\
Kombination über $v$ & $m_e/E_P$ & $-0{,}93\,\%$ & Summe der beiden vorigen \\
\bottomrule
\end{tabular}
\end{center}

Die Wege widersprechen sich nicht. Ihre Abweichungen sind je einer bekannten
Stelle zugeordnet: dem bare-Rest der Leptonleiter auf der einen, der
Identifikation von $E_\star$ auf der anderen Seite. Für die Vergleichstabelle
ist deshalb neben dem Wert immer der Weg anzugeben, auf dem er gewonnen wurde.

\paragraph{Folge für \#740.} Mit der eigenen FFGFT-Kette werden die
Brückengleichungen absolut. Die OPH-Zellkante ist
$L_{\rm cell}=\sqrt P\,\ell_P=2{,}064\times10^{-35}$\,m, der Zellentakt
$\tilde T_{\rm cell}=\sqrt P\,t_P=6{,}885\times10^{-44}$\,s \KM; jedes weitere
Taktverhältnis ist nach $\tilde T\cdot m=1$ ein Massenverhältnis. Länge und
Takt folgen damit ohne den Umweg über $N$ und $\Lambda$, unter der Lesart
$E_\star=E_P$. Deren Koeffizient ist auf OPH-Seite inzwischen abgeleitet
(Dok.~376): aus zwei Beobachtern mit verschiedener Spannung folgt
$\kappa=b\,L^2/q$ und damit $G=L^2$ \QM. \#740 selbst bleibt dort offen;
Längenidentität, Takt und ein Matching-Faktor werden weiterhin als offen
geführt. Die Skalenfrage aus \#740 ist damit auf FFGFT-Seite beantwortet, der
Punkt als ganzer nicht.

\section{Was die Brücke nicht leistet}

\begin{itemize}\itemsep2pt
\item Der obere Teil der Kette stammt aus einem Fremdrahmen und ist dort
  bedingt \QM; die Kombination ist ein Konsistenztest zwischen den Rahmen.
\item Zeiteinheit: In FFGFT folgt jede Zeitskala aus $\tilde T=1/m$; die
  SI-Sekunde dient nur der Umrechnung. In OPH ist die Zeiteinheit an $E_\star$
  gebunden und über die Uhrenanbindung festzulegen; dieser Schritt ist auf
  OPH-Seite offen \SM. Über die FFGFT-Kette ist der Zellentakt unter der
  Lesart $E_\star=E_P$ bereits absolut; die Laboruhr ist dann Prüfung.
\item $N$ und der kosmische Sektor sind nicht beteiligt (P39).
\item Die algebraischen Träger unterscheiden sich weiterhin ($A_5$ gegen
  $T^4/\mathbb Z_3$, Dok.~364); die Brücke betrifft Skalen, nicht Symmetrien.
\end{itemize}

\section*{Werkzeug und Methode}
\addcontentsline{toc}{section}{Werkzeug und Methode}
Claude (Anthropic) wurde als Werkzeug eingesetzt für: Lesen der
OPH-Quelldateien, Schreiben und Ausführen des Python-Prüfskripts, Ausarbeiten
der \LaTeX-Quellen und Pflege von Register und Changelog. Das Werkzeug rechnet
und schreibt; der Autor entscheidet, was berechnet und geschrieben wird,
beurteilt jedes Ergebnis und bestimmt, was es physikalisch bedeutet. Jede
Aussage ist durch Assertions im Prüfskript abgesichert. OPH-Werte sind wörtlich
aus dem öffentlichen Repositorium übernommen; Messwerte (CODATA 2022, PDG)
dienen nur als Komparator.

\section{Ergebnis}

\begin{enumerate}
\item B1: $E_{\rm cell}^{\rm OPH}=E_{\rm bit}(\sqrt P\,\ell_P)$ \BM.
\item B2: $n_{\rm thr}(L_{\rm cell})=\sqrt{P/2}$, Grenze $P=2$, OPH darunter \BM.
\item B3: Zellentakt $\tilde T_{\rm cell}=\sqrt P\,t_P$ aus $\tilde T\cdot m=1$;
  Laboranbindung auf OPH-Seite offen \SM.
\item B4: $L_{\rm cell}/L_0=\sqrt P/\xi\approx 9578$ \KM.
\item Treffpunkt $v$: OPH $v/E_\star$ \QM{} $\times$ FFGFT $m_e/v$ \KM{} ergibt
  $m_e/E_P$ auf $-0{,}93\,\%$; ein einziger SI-Anker genügt für beide Rahmen.
\item Länge und Takt absolut über die FFGFT-Kette: $L_{\rm cell}=2{,}064\times10^{-35}$\,m,
  $\tilde T_{\rm cell}=6{,}885\times10^{-44}$\,s \KM; Koeffizient der Lesart auf
  OPH-Seite abgeleitet, \#740 als ganzer dort offen (Dok.~376).
\item Eigene FFGFT-Kette $\xi\to G\to E_P$: $m_e/E_P$ auf $+0{,}002\,\%$ \KM;
  die Genauigkeit hängt vom Rechenweg ab (§\ref{sec:wege}).
\end{enumerate}

\medskip
\noindent\textbf{Prüfskript:} \texttt{2/python/Dok374\_Skripte/pruef\_374\_anker.py} (30/30)

\medskip
\noindent\textbf{Register (Dok.~190):} kein Eintrag.

\begin{thebibliography}{9}
\bibitem{dok006} J.~Pascher, \emph{Dok.~006 --- T0-Teilchenmassen} (2025).
\bibitem{dok012} J.~Pascher, \emph{Dok.~012 --- T0-Gravitationskonstante} (2025).
\bibitem{dok013} J.~Pascher, \emph{Dok.~013 --- T0 und SI-Einheiten} (2025).
\bibitem{dok180} J.~Pascher, \emph{Dok.~180 --- Herleitung von $L_0$} (2026).
\bibitem{dok257} J.~Pascher, \emph{Dok.~257 --- Bitenergie} (2026).
\bibitem{dok329} J.~Pascher, \emph{Dok.~329 --- $n_{\rm thr}$-Skalenanalyse} (2026).
\bibitem{dok352} J.~Pascher, \emph{Dok.~352 --- Leptonreste} (2026).
\bibitem{dok364} J.~Pascher, \emph{Dok.~364 --- Der Ikosaeder-Träger und das $D_4$-Gitter} (2026).
\bibitem{dok373} J.~Pascher, \emph{Dok.~373 --- Der Yukawa-Mechanismus} (2026).
\bibitem{oph} B.~Mueller et al., \emph{Observer Patch Holography},
  \texttt{github.com/FloatingPragma/observer-patch-holography} (Stand 22.~Sept.~2026).
\end{thebibliography}
